\scanpage{130}
se ramène à celle des hom composants.

À ce sujet, sous réserve de vérifications, je présume

$1^\circ)$ que l'hom $\mathcal E_{\tilde U_x}^{\bar L}\longrightarrow\mathcal E_U^{\bar L}$ est \uline{injectif}

--- ou ce qui revient au \unsure{=|même}, que sa partie géométrique
\[
\diagram{\begin{tikzcd}[column sep=large,row sep=small]
\mathcal E_{\tilde U_x/K}^{\bar L}\arrow[r]\arrow[d,equal] & \pi_{U/K}^{\bar L}\\
I_x\simeq T^{(I(x))} &
\end{tikzcd}}
\]
\margin{les deux lettres $I$ ne sont pas les mêmes}
est \uline{injective}.

\margin{C'est visiblement \uline{faux} \inserted{\uline{« en général »}} pour $\operatorname{card}I(x)\ne1$, comme on voit en \unsure{passant} à un autre modèle $X'$ déduit de $X$ en y faisant éclater un pt de $D$\ldots\ Mais y a-t-il un modèle « minimal » $X\supset U$, et pour un tel modèle \unsure{a-t-on} l'injectivité ???}

$2^\circ)$ que \deleted{l'hom} lorsque le corps résiduel de $V$ est alg $/K$ (i.e. $V$ correspondant à \unsure{un pt} de $\unsure{\mathfrak X_{L/K}|\mathcal Z_{L/K}}$ de dim zéro), alors l'hom
\[
 \mathcal E_{\tilde\eta}^{\bar L}=N_V\longrightarrow\mathcal E_{\tilde U_x}^{\bar L}
\]
\deleted{est injectif} \inserted{a une image} \deleted{surjectif} \inserted{ouverte} --- ce qui revient ici (puisque $k(V)$ \deleted{alg} \inserted{fini} sur $k(x)$ donc $\mathcal E_{k(V)}=N_V/I_V\hookrightarrow\mathcal E_{k(x)}$ \inserted{(sous-groupe ouv.)} $\simeq N_x/I_x$), que la partie géométrique
\[
 I_V\longrightarrow I_x=T^{I(x)}
\]
\deleted{est injective} \deleted{\unsure{surjective}} \inserted{a une image ouverte} \deleted{--- \unsure{je pense que} \illegible{3} \unsure{à un sous-groupe ouvert}.} \inserted{--- sans doute est-elle \unsure{=|même} surjective.}

À noter que \unsure{la} question $2^\circ$ est purement
\margin{129}
